% Generated draft: requires mathematical and source review.
\documentclass[11pt]{article}
\usepackage[margin=25mm]{geometry}
\usepackage{fontspec}
\setmainfont{Latin Modern Roman}
\usepackage{amsmath,amssymb,mathtools,unicode-math}
\setmathfont{Latin Modern Math}
\usepackage{enumitem,needspace,xurl,hyperref}
\setlength{\parindent}{0pt}
\setlength{\parskip}{5pt}
\setlength{\emergencystretch}{3em}
\newcommand{\N}{\mathbb{N}}
\newcommand{\Z}{\mathbb{Z}}
\newcommand{\Q}{\mathbb{Q}}
\newcommand{\R}{\mathbb{R}}
\newcommand{\C}{\mathbb{C}}
\newcommand{\F}{\mathbb{F}}
\newcommand{\A}{\mathbb{A}}
\newcommand{\PP}{\mathbb{P}}
\newcommand{\E}{\mathbb{E}}
\newcommand{\eps}{\varepsilon}
\newcommand{\dd}{\,\mathrm{d}}
\begin{document}
\section*{UnsolvedMath ID: 20002008}
\textbf{Source problem code:} catalogue entry
\subsection{Definitions and mathematical statement}
Section C.3 is a collection of seven computational questions, not one precisely specified algorithmic problem. Its items are: (A) compute examples of amoebas and tropical varieties; (B) compute homology of amoeba complements and seek a Hermitian form whose signature gives the number of complement components; (C) determine expansiveness or membership of zero for classical varieties, listing Grassmannians, generalized flag manifolds, spherical varieties, determinantal varieties, Schubert varieties and Calabi--Yaus; (D) develop a theory of smooth curves; (E) determine the complexity of resultant vanishing; (F) decide when \(\operatorname{Log}(X)\) lies in a lower-dimensional subspace; and (G) determine the complexity of distance to a discriminantal variety.

The aid makes item (F) precise in the following irreducible algebraic setting. Let
\[
 \varnothing\ne X\subset(\C^*)^n
\]
be a closed irreducible subvariety and let
\(\operatorname{Log}(z)=(\log|z_1|,\ldots,\log|z_n|)\).
Define its lattice of constant characters
\[
 L_X=\{u\in\Z^n:\exists c_u\in\C^*,\ z^u=c_u\text{ on }X\}.
\]
Equivalently, \(x^u-c_u\in I(X)\). Irreducibility makes \(L_X\) saturated: if a positive multiple of a character is constant, that character itself is constant. The map \(u\mapsto c_u\) is multiplicative. The aid's affine-hull theorem is
\[
 \operatorname{Aff}_{\R}(\operatorname{Log}(X))
 =\{a\in\R^n:\langle u,a\rangle=\log|c_u|
                         \text{ for every }u\in L_X\}.
\]
Its codimension is \(\operatorname{rank}L_X\). Every real affine equation on the amoeba is a real linear combination of these character equations. Thus containment in a proper affine subspace is equivalent to the existence of a nonzero constant integral character. The direction space is rational; its offsets \(\log|c_u|\) need not be rational numbers.

For \(x_0\in X\), the corresponding algebraic set
\[
 x_0T_X=\{z\in(\C^*)^n:z^u=c_u\ (u\in L_X)\}
\]
is a coset of a connected algebraic subtorus, by saturation, and is the smallest torus coset containing \(X\). This gives the stated logarithmic shadow of the minimal coset. The theorem concerns affine-hull dimension, not the local real dimension of the amoeba; a small-dimensional amoeba can have full affine span. Nor does it automatically extend to reducible \(X\) without taking the affine span of its component data.

If ``subspace'' in item (F) means a linear subspace through zero, there is a further distinction. Extend \(u\mapsto\log|c_u|\) linearly to \(\ell:L_X\otimes\R\to\R\). The homogeneous real equations are exactly
\[
 \langle v,a\rangle=0\quad\text{with }v\in L_X\otimes\R,
 \quad\ell(v)=0.
\]
A nonzero such \(v\), rather than merely a nonzero \(L_X\), is required for proper linear containment. The source does not specify affine versus through-origin meaning.

The other items remain open-ended in this record. For (C), the well-defined membership test in a fixed torus embedding is
\[
 0\in\operatorname{Log}(X)\quad\Longleftrightarrow\quad
 X\cap(S^1)^n\ne\varnothing.
\]
The word expansive commonly refers in the associated algebraic-dynamical convention to avoidance of this unit torus, whereas the source's parenthesis asks whether zero \emph{is} in the amoeba; the convention must be fixed. In every listed classical family the torus embedding and defining data matter, not just its abstract variety type. For (A)--(B), input class, coefficient field, homology coefficients and a nontrivial canonical Hermitian-form construction are unspecified. Item (D) does not state which algebraic or tropical smoothness notion is intended. For (E) and (G), coefficient encoding, dense/sparse presentation, variable/support bounds, computational model, norm, exact decision versus approximation and accuracy are absent. Therefore no definite complexity bound, random-model conclusion or classification of these six items is inferred from the affine-hull theorem.
\subsection{Short English statement}
For a nonempty irreducible torus subvariety, all affine equations on its amoeba come from constant integral characters. Their lattice determines the affine codimension and smallest containing torus coset. This resolves a precise affine version of item (F); the other computational questions still need their input and output conventions specified.
\subsection{Sources}
\begin{itemize}
\item Permanent UnsolvedMath ID \(20002008\), exact public code catalogue entry; local registry \texttt{aim-geometry-notes.json}, \texttt{notes[345]}, canonical tag section.
\item AIM workshop ``Amoebas and tropical geometry,'' section C.3 ``Computational issues'': \url{https://aimath.org/WWN/amoebas/amoebas.pdf}. No PDF was accessed. The terminal 17 is retained only as an extraction/page artifact, not an item of the mathematics.
\item Research attempt 1, ``Affine dimension of algebraic amoebas via constant characters,'' labelled partial\_result, partially\_solved, proved and low novelty confidence. It mentions a literature-based synthesis for other items but supplies no detailed synthesis or citations in this record. No proposer, year or original literature reference is supplied.
\end{itemize}
\end{document}
